%% build: latexmk -pdf -pvc -xelatex prez.tex

\documentclass[dvipsnames,aspectratio=169]{beamer}
\usetheme{Madrid}


%% kill footline
\setbeamertemplate{footline}[frame number]{}
\setbeamertemplate{navigation symbols}{}
\setbeamertemplate{footline}{}

%% bibliography
\bibliographystyle{alpha}
\setbeamerfont{bibliography item}{size=\footnotesize}
\setbeamerfont{bibliography entry author}{size=\footnotesize}
\setbeamerfont{bibliography entry title}{size=\footnotesize}
\setbeamerfont{bibliography entry location}{size=\footnotesize}
\setbeamerfont{bibliography entry note}{size=\footnotesize}
\setbeamertemplate{bibliography item}{}

%% kill ball enumeration
\setbeamertemplate{enumerate items}[circle]
\setbeamertemplate{section in toc}[circle]

%% kill block shadows
\setbeamertemplate{blocks}[rounded][shadow=false]
\setbeamertemplate{title page}[default][colsep=-4bp,rounded=true]

%% kill ball itemize
\setbeamertemplate{itemize items}[circle]

%% --------------------------------------------------------------------------------

\usepackage[utf8]{inputenc}
\usepackage{amsmath}
\usepackage{cite}
\usepackage{amsthm}
\usepackage{amssymb}
\usepackage{amsfonts}
\usepackage{mathpartir}
\usepackage{scalerel}
\usepackage{stmaryrd}
\usepackage{bm}
\usepackage{graphicx}
\usepackage{fontspec}
\usepackage{textcomp}
\usepackage{fancyvrb}
\usepackage{tikz-cd}

\tikzcdset{scale cd/.style={every label/.append style={scale=#1},
    cells={nodes={scale=#1}}}}

\setmonofont[Scale=0.8]{DejaVu Sans Mono Bold}

%% --------------------------------------------------------------------------------

%% HoTT style composition
\makeatletter
\DeclareRobustCommand{\sqcdot}{\mathbin{\mathpalette\morphic@sqcdot\relax}}
\newcommand{\morphic@sqcdot}[2]{%
  \sbox\z@{$\m@th#1\centerdot$}%
  \ht\z@=.33333\ht\z@
  \vcenter{\box\z@}%
}
\makeatother

\renewcommand{\mit}[1]{{\mathit{#1}}}
\newcommand{\ttt}[1]{{\texttt{#1}}}
\newcommand{\msf}[1]{{\mathsf{#1}}}
\newcommand{\mbf}[1]{{\mathbf{#1}}}
\newcommand{\mbb}[1]{\mathbb{#1}}
\newcommand{\vsp}{\vspace{1em}}
\newcommand{\vsph}{\vspace{0.5em}}

%% \newcommand{\U}{\mathsf{U}}
%% \newcommand{\code}{\mathsf{code}}
%% \newcommand{\bs}[1]{\boldsymbol{#1}}
%% \newcommand{\wh}[1]{\widehat{#1}}
%% \newcommand{\mdo}{\mbf{do}\,}
%% \newcommand{\ind}{\hspace{1em}}
%% \newcommand{\bif}{\mbf{if}\,}
%% \newcommand{\bthen}{\mbf{then}\,}
%% \newcommand{\belse}{\mbf{else}\,}
%% \newcommand{\return}{\mbf{return}\,}
%% \newcommand{\pure}{\mbf{pure}\,}
%% \newcommand{\lam}{\lambda\,}
%% \newcommand{\data}{\mbf{data}\,}
%% \newcommand{\where}{\mbf{where}}
%% \newcommand{\M}{\msf{M}}
%% \newcommand{\letrec}{\mbf{letrec}\,}
%% \newcommand{\of}{\mbf{of}\,}
%% \newcommand{\go}{\mit{go}}
%% \newcommand{\add}{\mit{add}}
%% \newcommand{\letdef}{\mbf{let\,}}
%% \newcommand{\map}{\mit{map}}
%% \newcommand{\emptycon}{\scaleobj{.75}\bullet}
%% \newcommand{\Tyo}{\msf{Ty}_{\mbbo}}
%% \newcommand{\Tmo}{\msf{Tm}_{\mbbo}}
%% \newcommand{\whset}{\wh{\Set}}
%% \newcommand{\ev}{\mbb{E}}
%% \newcommand{\re}{\mbb{R}}


%% \newcommand{\mbbc}{\mbb{C}}
%% \newcommand{\mbbo}{\mbb{O}}

%% \newcommand{\vas}{\mathsf{as}}
%% \newcommand{\vbs}{\mathsf{bs}}
%% \newcommand{\vcs}{\mathsf{cs}}
%% \newcommand{\vxs}{\mathsf{xs}}
%% \newcommand{\vys}{\mathsf{ys}}
%% \newcommand{\vsp}{\mathsf{sp}}
%% \newcommand{\vma}{\mathsf{ma}}
%% \newcommand{\vga}{\mathsf{ga}}
%% \newcommand{\vm}{\mathsf{m}}
%% \newcommand{\vn}{\mathsf{n}}
%% \newcommand{\vk}{\mathsf{k}}
%% \newcommand{\vA}{\mathsf{A}}
%% \newcommand{\vB}{\mathsf{B}}
%% \newcommand{\vC}{\mathsf{C}}
%% \newcommand{\vS}{\mathsf{S}}
%% \newcommand{\vF}{\mathsf{F}}
%% \newcommand{\vR}{\mathsf{R}}
%% \newcommand{\vM}{\mathsf{M}}
%% \newcommand{\vmb}{\mathsf{mb}}
%% \newcommand{\mAs}{\mathsf{As}}
%% \newcommand{\va}{\mathsf{a}}
%% \newcommand{\vb}{\mathsf{b}}
%% \newcommand{\vc}{\mathsf{c}}
%% \newcommand{\vd}{\mathsf{d}}
%% \newcommand{\vx}{\mathsf{x}}
%% \newcommand{\vy}{\mathsf{y}}
%% \newcommand{\vz}{\mathsf{z}}
%% \newcommand{\vf}{\mathsf{f}}
%% \newcommand{\vfs}{\mathsf{fs}}
%% \newcommand{\vg}{\mathsf{g}}
%% \newcommand{\vh}{\mathsf{h}}
%% \newcommand{\vt}{\mathsf{t}}
%% \newcommand{\vs}{\mathsf{s}}
%% \newcommand{\vr}{\mathsf{r}}
%% \newcommand{\vu}{\mathsf{u}}
%% \newcommand{\vl}{\mathsf{l}}
%% \newcommand{\vns}{\mathsf{ns}}
%% \newcommand{\vW}{\mathsf{W}}
%% \newcommand{\vsup}{\mathsf{sup}}
%% \newcommand{\vid}{\mathsf{id}}
%% \newcommand{\whW}{\wh{\vW}}


%% \newcommand{\SOP}{\msf{SOP}}
%% \newcommand{\El}{\msf{El}}
%% \newcommand{\USOP}{\msf{U}_{\msf{SOP}}}
%% \newcommand{\Uprod}{\msf{U_P}}
%% \newcommand{\Elprod}{\msf{El_{P}}}
%% \newcommand{\IsSOP}{\msf{IsSOP}}
%% \newcommand{\forEach}{\msf{forEach}}
%% \newcommand{\single}{\msf{single}}
%% \newcommand{\msplit}{\msf{split}}
%% \newcommand{\mapGen}{\msf{mapGen}}
%% \newcommand{\genPull}{\msf{gen_{Pull}}}
%% \newcommand{\Set}{\msf{Set}}
%% \newcommand{\casePull}{\msf{case_{Pull}}}
%% \newcommand{\appull}{\ap_{\Pull}}

%% \newcommand{\Con}{\msf{Con}}
%% \newcommand{\Sub}{\msf{Sub}}
%% \newcommand{\Tm}{\msf{Tm}}
%% \newcommand{\Mor}{\msf{Tm}}

%% \newcommand{\ext}{\triangleright}

%% \newcommand{\Int}{\msf{Int}}
%% \newcommand{\List}{\msf{List}}
%% \newcommand{\Tree}{\msf{Tree}}
%% \newcommand{\Node}{\msf{Node}}
%% \newcommand{\Leaf}{\msf{Leaf}}
%% \newcommand{\Nil}{\msf{Nil}}
%% \newcommand{\Cons}{\msf{Cons}}
%% \newcommand{\Reader}{\msf{Reader}}
%% \newcommand{\ReaderT}{\msf{ReaderT}}
%% \newcommand{\Monad}{\msf{Monad}}
%% \newcommand{\Applicative}{\msf{Applicative}}
%% \newcommand{\class}{\msf{class}}
%% \newcommand{\Functor}{\msf{Functor}}
%% \newcommand{\Bool}{\msf{Bool}}
%% \newcommand{\Statel}{\msf{State}}
%% \newcommand{\fro}{\leftarrow}
%% \newcommand{\case}{\mbf{case\,}}
%% \newcommand{\foldr}{\msf{foldr}}
%% \newcommand{\foldl}{\msf{foldl}}
%% \newcommand{\rep}{\msf{rep}}
%% \newcommand{\concatMap}{\msf{concatMap}}

%% \newcommand{\Lift}{{\Uparrow}}
%% \newcommand{\Up}{{\Uparrow}}
%% \newcommand{\spl}{{\bs{\sim}}}
%% \newcommand{\ql}{{\bs{\langle}}}
%% \newcommand{\qr}{{\bs{\rangle}}}
%% \newcommand{\bind}{\mathbin{>\!\!>\mkern-6.7mu=}}

%% \newcommand{\MTy}{\msf{MetaTy}}
%% \newcommand{\MTm}{\msf{MetaTm}}
%% \newcommand{\VTy}{\msf{ValTy}}
%% \newcommand{\Ty}{\msf{Ty}}
%% \newcommand{\CTy}{\msf{CompTy}}
%% \newcommand{\True}{\msf{True}}
%% \newcommand{\False}{\msf{False}}
%% \newcommand{\fst}{\msf{fst}}
%% \newcommand{\snd}{\msf{snd}}

\newcommand{\blank}{{\mathord{\hspace{1pt}\text{--}\hspace{1pt}}}}

%% \newcommand{\Nat}{\msf{Nat}}
%% \newcommand{\Zero}{\msf{Zero}}
%% \newcommand{\Suc}{\msf{Suc}}
%% \newcommand{\Maybe}{\msf{Maybe}}
%% \newcommand{\MaybeT}{\msf{MaybeT}}
%% \newcommand{\Nothing}{\msf{Nothing}}
%% \newcommand{\Just}{\msf{Just}}

\theoremstyle{remark}
\newtheorem{notation}{Notation}
\newtheorem*{axiom}{Axiom}

%% \newcommand{\id}{\mit{id}}
%% \newcommand{\mup}{\mbf{up}}
%% \newcommand{\mdown}{\mbf{down}}
%% \newcommand{\tyclass}{\mbf{class}}
%% \newcommand{\instance}{\mbf{instance}\,}
%% \newcommand{\Improve}{\msf{Improve}}
%% \newcommand{\Gen}{\msf{Gen}}
%% \newcommand{\unGen}{\mit{unGen}}
%% \renewcommand{\Vec}{\msf{Vec}}
%% \newcommand{\gen}{\mit{gen}}
%% \newcommand{\genRec}{\mit{genRec}}
%% \newcommand{\fmap}{<\!\!\$\!\!>}
%% \newcommand{\ap}
%% \newcommand{\runGen}{\mit{runGen}}
%% \newcommand{\qt}[1]{\ql#1\qr}
%% \newcommand{\lift}{\mit{lift}}
%% \newcommand{\liftGen}{\mit{liftGen}}
%% \newcommand{\MonadGen}{\msf{MonadGen}}
%% \newcommand{\MonadState}{\msf{MonadState}}
%% \newcommand{\MonadReader}{\msf{MonadReader}}
%% \newcommand{\RA}{\Rightarrow}
%% \newcommand{\EitherT}{\msf{EitherT}}
%% \newcommand{\Either}{\msf{Either}}
%% \newcommand{\Left}{\msf{Left}}
%% \newcommand{\Right}{\msf{Right}}
%% \newcommand{\StateT}{\msf{StateT}}
%% \newcommand{\Identity}{\msf{Identity}}

%% \newcommand{\Stop}{\msf{Stop}}
%% \newcommand{\Skip}{\msf{Skip}}
%% \newcommand{\Yield}{\msf{Yield}}

%% \newcommand{\runIdentity}{\mit{runIdentity}}
%% \newcommand{\runReaderT}{\mit{runReaderT}}
%% \newcommand{\newtype}{\mbf{newtype}\,}
%% \newcommand{\runMaybeT}{\mit{runMaybeT}}
%% \newcommand{\runStateT}{\mit{runStateT}}
%% \newcommand{\runState}{\mit{runState}}
%% \newcommand{\dlr}{\,\$\,}
%% \newcommand{\ImproveF}{\msf{ImproveF}}
%% \newcommand{\ExceptT}{\msf{ExceptT}}
%% \newcommand{\State}{\msf{State}}
%% \newcommand{\SumVS}{\msf{SumVS}}
%% \newcommand{\ProdCS}{\msf{ProdCS}}
%% \newcommand{\Here}{\msf{Here}}
%% \newcommand{\There}{\msf{There}}
%% \newcommand{\IsSumVS}{\msf{IsSumVS}}
%% \newcommand{\MonadJoin}{\msf{MonadJoin}}
%% \newcommand{\Stream}{\msf{Stream}}
%% \newcommand{\join}{\mit{join}}
%% \newcommand{\modify}{\mit{modify}}
%% \newcommand{\get}{\mit{get}}
%% \newcommand{\mput}{\mit{put}}
%% \newcommand{\Rep}{\mit{Rep}}
%% \newcommand{\encode}{\mit{encode}}
%% \newcommand{\decode}{\mit{decode}}
%% \newcommand{\mindex}{\mit{index}}
%% \newcommand{\mtabulate}{\mit{tabulate}}
%% \newcommand{\States}{\mit{States}}
%% \newcommand{\seed}{\mit{seed}}
%% \newcommand{\step}{\mit{step}}
%% \newcommand{\Step}{\msf{Step}}
%% \newcommand{\Pull}{\msf{Pull}}
%% \newcommand{\MkPull}{\msf{MkPull}}

\newcommand{\RED}[1]{{\color{BrickRed} #1}}
\newcommand{\GRN}[1]{{\color{OliveGreen} #1}}
\newcommand{\RBL}[1]{{\color{RoyalBlue} #1}}
\newcommand{\BLU}[1]{{\color{Blue} #1}}

%% \newcommand{\bsl}{\textbackslash}
%% \newcommand{\brl}{\textbraceleft}
%% \newcommand{\brr}{\textbraceright}

%% --------------------------------------------------------------------------------

\title{Thinning-Stable Metaprogramming}
\author{András Kovács}
\institute{University of Cambridge}
\date{13rd Oct 2026, IFIP WG 2.11 meeting, Niagara-on-the-Lake}
\begin{document}

\frame{\titlepage}

\begin{frame}[fragile]{Stability}

\vsp
In metaprogramming, it's nice if we don't have to explicitly handle
object-language scopes.
\vsp

\begin{Verbatim}[commandchars=\\\{\}]
  f : Exp A -> Exp A
\end{Verbatim}
\vsp
We'd like \texttt{f} to work on expressions containing arbitrary free variables.
\vsp

At least, we'd like a closed such \texttt{f} to have the following semantics
\vsp
\begin{Verbatim}[commandchars=\\\{\}]
  ⟦f⟧ : ∀ Γ. Exp Γ A -> Exp Γ A
\end{Verbatim}
\vsp

\end{frame}

\begin{frame}[fragile]{Stability}

Just having \texttt{⟦f⟧ : ∀ Γ. Exp Γ A -> Exp Γ A} is not enough.
\vsp


\textbf{First}, there's no \emph{hygiene}. \texttt{f} could inspect \texttt{Γ} in arbitrary ways,
and e.g.\ return different expressions in empty and non-empty scopes.

\end{frame}


\begin{frame}[fragile]{Stability}

\textbf{Second}, There's not even an obvious way to have \emph{binders}. With STLC as an object language,
we'd like at least
\vsph
\begin{Verbatim}[commandchars=\\\{\}]
  lam : (Exp A -> Exp B) -> Exp (A ⇒ B)
\end{Verbatim}
\vsph
However, the standard indexed-set semantics demands
\vsph
\begin{Verbatim}[commandchars=\\\{\}]
  ⟦lam⟧ : ∀ Γ. (Exp Γ A -> Exp Γ B) -> Exp Γ (A ⇒ B)
\end{Verbatim}
\vsph
where we can't conjure a fresh variable to apply the inner function to.

\end{frame}


\begin{frame}[fragile]{Explicit contexts}

One approach: have explicit contexts in user-facing language.
\vsp
\begin{Verbatim}[commandchars=\\\{\}]
  f : ∀ Γ. Exp Γ A → Exp Γ A
\end{Verbatim}
\vsp
For example: Moebius, Cocon, layered modal TT, typed FOAS.
\vsp

\begin{itemize}
\item We can do \emph{arbitrary} intensional analysis.
\item A lot of bureaucracy; ergonomics is a challenge.
\end{itemize}

\end{frame}


\begin{frame}[fragile]{Stability under substitution}

\begin{Verbatim}[commandchars=\\\{\}]
  ⟦f⟧ : ∀ Γ. Exp Γ A → Exp Γ A
    such that
  ∀ (Δ : Scope) (σ : Sub Δ Γ) (a : Exp Γ A). ⟦f⟧ (a[σ]) = (⟦f⟧ a)[σ]
\end{Verbatim}
\vsp
A standard semantics for HOAS, logical frameworks \& two-level calculi.
\vsp

\emph{Barely any intensional analysis}: no expression can be distinguished from
a variable.

\vsp
Easy faithful representation of object binders:
\begin{Verbatim}[commandchars=\\\{\}]
  lam   : (Exp A -> Exp B) -> Exp (A ⇒ B)
  ⟦lam⟧ : ∀ Γ. (∀ Δ. Sub Δ Γ -> Exp Δ A → Exp Δ B) -> Exp Γ (A ⇒ B)
\end{Verbatim}

\end{frame}

\begin{frame}[fragile]{Stability under variable renaming}

A renaming from \texttt{Γ} to \texttt{Δ} maps every \texttt{Δ}-variable to
a \texttt{Γ}-variable.
\vsp

To faithfully represent binders, we need an extra \emph{type of variables}, and e.g.\ have
\vsp
\begin{Verbatim}[commandchars=\\\{\}]
  lam : (Var A -> Exp B) -> Exp (A ⇒ B)
\end{Verbatim}
\vsp

Metaprograms can distinguish variables and non-variables, but e.g.\ can't decide equality
of expressions...

\vsp
... because two distinct variables can become equal by renaming.

\vsp
Used in metatheory of TTs, but rarely in metaprogramming (?)

\end{frame}

\begin{frame}[fragile]{Stability under injective variable renaming}

Equality of expressions can be decided.
\vsp

But binders are much more fiddly: it's a \emph{sub-structural} function type.
\vsp

Example: \emph{nominal sets}.
\vsp

\textbf{Example:} the nullary parametric type theory (NPTT) of Muylder, Nuyts and Devriese\footnote{Muylder, Nuyts, Devriese: \emph{Nominal Type Theory by Nullary Internal Parametricity}}.

\end{frame}


\begin{frame}[fragile]{Stability under thinning}

A thinning from \texttt{Γ} to \texttt{Δ} deletes zero of more binders from \texttt{Γ},
resulting in \texttt{Δ}.
\vsp

Contravariant action: if \texttt{σ : Thinning Γ Δ} and \texttt{e : Exp Δ A},
then \texttt{e[σ] : Exp Γ A}.

\vsp It's the \emph{weakest} stability requirement that makes sense for hygienic
metaprogramming.

\end{frame}


\begin{frame}[fragile]{ThinLF}

A mostly hypothetical language. Agda/Idris syntax with extra nominal
features. Based on NPTT by Muylder, Nuyts \& Devriese.

\vsp

\begin{Verbatim}[commandchars=\\\{\}]
  data Ty : Type where
    ι   : Ty
    _⇒_ : Ty → Ty → Ty

  data Exp : Ty → NomType where
    var : Var (Exp A) → Exp A
    lam : (@ (Exp A) → Exp B) → Exp (A ⇒ B)
    app : Exp (A ⇒ B) → Exp A → Exp B
\end{Verbatim} 

\end{frame}

\begin{frame}[fragile]{ThinLF}
We can declare an inductive type to be nominal, and use \texttt{Var} and
the nominal function type in the definition.
\vsp
\begin{Verbatim}[commandchars=\\\{\}]
    Γ context   Γ ⊢ A : NomType     Γ ⊢ A : NomType    Γ, x @ A ⊢ B : Type
    ───────────────────────────     ──────────────────────────────────────
         Γ, x @ A context                  Γ ⊢ (x @ A) → B : Type


              Γ, x @ A ⊢ t : B           Γ ⊢ t : (x @ A) → B
           ────────────────────────    ────────────────────────
           Γ ⊢ λ x. t : (x @ A) → B    Γ,y @ A,Δ ⊢ t y : B[y/x]

             Γ ⊢ A : NomType
             ────────────────    ────────────────────────
             Γ ⊢ Var A : Type    Γ, x @ A, Δ ⊢ #x : Var A
\end{Verbatim}
\vsp

\end{frame}

\begin{frame}[fragile]{ThinLF}

Ordinary HOAS/LF programming is still available:
\begin{Verbatim}[commandchars=\\\{\}]
  lam' : (Exp A → Exp B) → Exp (A ⇒ B)
  lam' f = lam λ x. f (var #x)
\end{Verbatim}
\vsp
No flipping:
\begin{Verbatim}[commandchars=\\\{\}]
  flip : (@ A → @ B → C) → (@ B → @ A → C)
  flip = λ f b a. f a b   \RED{-- ill-typed}
\end{Verbatim}
\vsp
No duplication:
\begin{Verbatim}[commandchars=\\\{\}]
  dup : (@ A → @ A → B) → @ A → B
  dup = λ f a. f a a   \RED{-- ill-typed}
\end{Verbatim}
\end{frame}

\begin{frame}[fragile]{Freshness}

\texttt{Fresh A x}: the type of \texttt{A}-values that
don't depend on \texttt{x}.
\vsp

\begin{Verbatim}[commandchars=\\\{\}]
  Fresh : \{A : NomType\} → Type → @ A → Type
  fresh : ((x @ A) → Fresh B x) → B
  <\_,\_> : B → (x @ A) → Fresh B x
  \GRN{ -- fresh and <\_,\_> are definitional inverses}
\end{Verbatim}
\vsp

We can match under a binder to extract a \texttt{Bool}:
\begin{Verbatim}[commandchars=\\\{\}]
  extractBool : (@ A → Bool) → Bool
  extractBool f = fresh λ x. case f x of
    True  → <True, x>
    False → <False, x>
\end{Verbatim} 
\end{frame}

\begin{frame}[fragile]{Abstracting over free variables}

Another primitive operation:
\begin{Verbatim}[commandchars=\\\{\}]
  rebind : ((@ A → B) → @ A → C) → (@ A → B → C)
\end{Verbatim}
(I only show the non-dependently typed version here).
\vsp

We can use it to match on \texttt{Either} under a binder:
\vsp
\begin{Verbatim}[commandchars=\\\{\}]
  split : (@ A → Either B C) → Either (@ A → B) (@ A → C)
  split f = fresh λ x. case f x of
    Left b  → rebind (λ b x. <Left b, x>) x b
    Right c → rebind (λ c x. <Right c, x>) x c
\end{Verbatim}
\vsp

\end{frame}
\begin{frame}[fragile]{Some useful definitions}

\begin{Verbatim}[commandchars=\\\{\}]
  forget : \{A : NomType\}\{B : Type\} → (x @ A) → Fresh B x → B
  forget = rebind (λ fr _. fresh fr)

  map : (B → C) → (x @ A) → Fresh B x → Fresh C x
  map f = rebind (λ b x. <f (fresh b), x>)

  bind : (x @ A) → Fresh B x → (B → Fresh C x) → Fresh C x
  bind x b k = k (forget x b)

  Fresh' : \{A : NomType\}\{B : Type\} → B → @ A → Type
  Fresh' b x = (b' : Fresh B x) × (forget b' = b)
\end{Verbatim}
\end{frame}

\begin{frame}[fragile]{Structure abstraction principles (SAP)}

Most SAP-s from Muylder et al.\ work fine:
\begin{Verbatim}[commandchars=\\\{\}]
  (@ A → ⊤)             ≃ ⊤
  (@ A → ⊥)             ≃ ⊥
  (@ A → (B + C))       ≃ ((@ A → B) + (@ A → C))
  (@ A → (B × C))       ≃ ((@ A → B) × (@ A → C))
  (@ A → B → C)         ≃ ((@ A → B) → (@ A → C))
  (@ A → List B)        ≃ List (@ A → B)
  ((x @ A) → f x = g x) ≃ (f = g)
  ((x @ A) → Fresh B x) ≃ B
\end{Verbatim}

\vsp
This one requires name swapping, so we need to make it primitive (or generalize \texttt{fresh}):

\begin{Verbatim}[commandchars=\\\{\}]
  ((x @ A) → Fresh (B x) y) ≃ Fresh ((x @ A) → B x) y
\end{Verbatim}
\end{frame}

\begin{frame}[fragile]{Example: counting variables}

\begin{Verbatim}[commandchars=\\\{\}]
  SAP-Nat : (@ A → Nat) → Nat
  SAP-Nat f = fresh λ x.
     let go : Nat → Fresh Nat x
         go Zero    = <Zero, x>
         go (Suc n) = map Suc x (go n)
     go (f x)

  countVars : Exp A → Nat
  countVars (var _)   = 1
  countVars (app t u) = countVars t + countVars u
  countVars (lam f)   = SAP-Nat (λ x. countVars (f x))
\end{Verbatim}

\end{frame}

\begin{frame}[fragile]{Example: recovering scoped syntax}

\begin{Verbatim}[commandchars=\\\{\}]
  Scope : Type
  Scope = List ((A : Ty) × Var (Exp A))

  inScope : (Γ : Scope) → \{A : Ty\} → Exp A → Type
  inScope Γ (var \{A\} x) = (A, x) ∈ Γ
  inScope Γ (app t u)   = inScope Γ t × inScope Γ u
  inScope Γ (lam f)     = (x @ Exp A) → inScope (Cons (A, #x) Γ) (f x)
\end{Verbatim}
\vsp

Note: \texttt{inScope} is decidable if we have decidable equality for \texttt{Var}.

\end{frame}


\begin{frame}[fragile]{Example: well-typed interpreter on scoped syntax}

\begin{Verbatim}[commandchars=\\\{\}]
  ⟦_⟧ : Ty → Type
  ⟦ι⟧     = Nat
  ⟦A ⇒ B⟧ = ⟦A⟧ → ⟦B⟧

  SAP-⟦⟧-Ty : (@ A → ⟦B⟧) → ⟦B⟧
  SAP-⟦⟧-Ty = ...

  ⟦_⟧ : Scope → Type
  ⟦Nil          ⟧ = ⊤
  ⟦Cons (A, _) Γ⟧ = ⟦A⟧ × ⟦Γ⟧

  ⟦_⟧ : \{Γ A\} → (t : Exp A) → inScope Γ t → ⟦Γ⟧ → ⟦A⟧
  ⟦var x  ⟧ xsc        γ = lookup x xsc γ
  ⟦app t u⟧ (tsc, usc) γ = (⟦t⟧ tsc γ) (⟦u⟧ usc γ)
  ⟦lam f  ⟧ fsc        γ = SAP-⟦⟧-Ty (λ x α. ⟦ f x ⟧ (fsc x) (α, γ))
\end{Verbatim}

\end{frame}

\begin{frame}[fragile]{Extra assumptions on variables}

The semantics supports
\begin{itemize}
  \item Decidable ordering on \texttt{Var}.
  \item Decidable comparison of a raw variable and a \texttt{Var}.
\end{itemize}
\vsp

If we have these, we can define
\begin{itemize}
 \item Decidable ordering (and equality) for \texttt{Exp}.
 \item General variable renaming as a recursive operation on \texttt{Exp}.
 \item Substitution as a recursive operation on \texttt{Exp}.
 \item Strengthening on \texttt{Exp} w.r.t.\ a raw variable, returning
       \texttt{Decide (Fresh (Exp A) x)}.
\end{itemize}
\vsp

I'm not sure about the exact specification of variable operations.

\end{frame}

\begin{frame}[fragile]{ThinLF compared to Nullary Parametric TT}

ThinLF has a lot of overlap with NPTT. What are the
differences?
\vsp

\begin{itemize}
\item ThinLF supports \emph{indexed nominal datatypes}, where the indexing types are ordinary ThinLF
  types. NPTT only supports untyped nominal types in a straightforward way, and trying to have
  first-class indexing types seems rather complicated.
\item NPTT supports swapping raw variables in instantiations, ThinLF doesn't.
\item ThinLF supports decidable ordering of variables, NPTT is only compatible with decidable equality.
\end{itemize}


\end{frame}

\begin{frame}[fragile]{Conclusions}

\begin{itemize}
\item We can do a lot in ThinLF, but I don't necessarily know how to do it efficiently and concisely.
\item I'm hoping that we can build nice machinery on top of the primitives.
\item Features like \emph{pattern matching under binders} could be elaborated to ThinLF.
\item I'm not quite sure about the set of primitives, or even the basic setup. Can we use ``observational'' binder types instead of manual SAP plugging?






\end{itemize}


\end{frame}


\end{document}
